\part{A reader's map of the implementation}
\label{part:source-map}

\chapter{From a working session to the file that owns it}

Read this part with the source catalogue open. Each linked path leads to a
complete listing later in this book; the numbers in its left margin are
source lines, starting again at one for each file. Wrapped lines have no
new number. The catalogue is a reference, not a demand to read every file
in alphabetical order. First follow a request through the host; then follow
one Scheme wrapper back through the same boundary.

The two books share this map and the same source catalogue. The earlier
chapters of the design book explain the decisions; those of the expert
guide teach the language and operating habits. This part connects both to
the implementation. The catalogue prints every tracked, non-test Go file
and every tracked \file{.scm} file, including the project init file and the
exfiltration scenarios. Go test files, including their embedded Scheme
examples, remain in the checkout; the reading routes below name the tests
to study. JavaScript, CSS, shell scripts and third-party dependencies are
not reproduced. The console's Go bridge is included, but understanding its
rendering also requires \file{internal/frontend/web/assets/}.

\section{The ownership map}

\begin{longtable}{@{}p{3.8cm}p{10.3cm}@{}}
\toprule
Owner & Question it answers \\
\midrule\endhead
\file{cmd/} & Which mode starts, which configuration is used, and which frontend owns interaction? \\
\file{internal/core/} & How is one session assembled, serialised, resumed and closed? \\
\file{internal/loop/} & Which Scheme strategy runs, and which host invariants constrain model requests and tool replies? \\
\file{internal/scheme/} & What does evaluation mean, and how does every primitive reach the guard? \\
\file{internal/prim/} & What does an allowed filesystem, process, git or integration operation actually do? \\
\file{internal/guard/} & May this declared operation run with these arguments and this origin? \\
\file{internal/window/} & Which logged events become the next provider-neutral request? \\
\file{internal/llm/} & How does that request become one provider's wire format? \\
\file{internal/log/} & What persists, and what can replay reconstruct? \\
\file{internal/frontend/} & How do people submit work, observe events and answer prompts? \\
\bottomrule
\end{longtable}

Do not start by reading all of \file{core/}. It connects almost everything.
Start with the types in \src{internal/core/session.go}, the construction
sequence in \src{internal/core/construct.go}, and one public operation.
When a method delegates to another package, read that package's contract
before its implementation. Returning to the caller then reveals which
state belongs to the session and which state belongs to the subsystem.

\section{Route one: starting a conversation}

Begin at \src{cmd/nodwyn/main.go} and \src{cmd/nodwyn/cli.go}. Follow the
chosen mode into \src{cmd/nodwyn/chat.go}, \src{cmd/nodwyn/tui.go} or
\src{cmd/nodwyn/daemon.go}; they are alternative ways to reach the host,
not separate implementations of the agent's policy.

In \src{internal/core/construct.go}, read \code{Options}, \code{New},
\code{initialize}, \code{newGuard} and \code{loadInit}. Follow their actual
calls rather than assuming the order of functions in the file is execution
order. A requested worktree is created before runtime initialisation,
because it changes the root from which project policy is read. Failure
paths matter: an unstarted session must not leave a misleading empty log
or an abandoned newly created worktree.

The interpreter loads its prelude; the host registers its primitives; the
loop installs its default Scheme; the project's init file runs; then the
session fixes its prompts and tool catalogue. In
\src{internal/loop/loop.go}, compare \code{New} with \code{Freeze}. A
function that can be looked up again on the next turn is different from a
tool schema already advertised to a model. This distinction explains most
questions about why a reload did or did not change a running session.

For configuration follow \src{internal/config/config.go}.
Model choice during a session belongs to
\src{internal/core/model.go}, and context sizing to
\src{internal/core/windowsize.go}. Use the adjacent configuration tests
and \file{internal/core/model\_switch\_test.go} to distinguish startup
defaults from a conversation's later model choice.

\section{Route two: one request and one filesystem read}

\begin{enumerate}
\item In \src{internal/core/session.go}, trace \code{Send} through the
  turn machinery. Find where the user message is recorded, where a busy
  session refuses overlapping work, and where cancellation is installed.
\item In \src{internal/loop/loop.go}, find the lookup of
  \code{nodwyn-turn}. Follow that into
  \src{internal/loop/strategy.scm}. The strategy chooses when to ask again;
  \src{internal/loop/model.go} owns making a model request.
\item In \src{internal/window/window.go}, follow \code{Build}. The next
  context is derived from events; it is not a mutable transcript owned by
  the provider. \src{internal/llm/llm.go} defines the shared request and
  reply types. Choose one adapter, such as \src{internal/llm/ollama.go},
  before comparing it with \src{internal/llm/responses.go}.
\item Return to the reply's tool calls. Scheme's \code{run-calls} can
  wrap dispatch; the host implementation in \src{internal/loop/tools.go}
  records the call, converts named JSON arguments and applies the primitive.
  \src{internal/loop/turn.go} tracks the calls that still need replies.
\item Find \code{applyPrim} in \src{internal/scheme/prim.go}. Read the
  guard call before reading any filesystem implementation. Follow it to
  \src{internal/guard/guard.go} and the path checkers in
  \src{internal/guard/checkers.go}. Only then read the \code{fs-read}
  registration and handler in \src{internal/prim/fs.go}.
\item Follow the returned value back: provenance is attached, the result
  is rendered and recorded, and the next \code{Build} sees it. Compare
  \src{internal/core/events.go}, \src{internal/core/history.go} and
  \src{internal/log/log.go}: frontend notification, in-memory history and
  durable append are related, but are not interchangeable sinks.
\end{enumerate}

To check your understanding, explain why a Scheme tool that calls
\code{fs-read} cannot bypass a denied path. The answer must mention the
inner primitive's metadata and \code{applyPrim}, not the wrapper's friendly
name. Read \file{internal/core/strategy\_test.go}, especially the tests
for fabricated calls, rewritten arguments and unanswered calls. The loop
must keep a valid transcript even when a strategy goes wrong.

\section{Route three: a credential-bearing request}

Read \src{internal/binding/binding.go} for the facts that authorise a
credential's use, then \src{internal/keychain/store.go} for storage and
\src{internal/keychain/acquire.go} for acquisition. A handle in Scheme is
not the key's text. The composition happens in
\src{internal/core/credentials.go} and \src{internal/prim/egress.go}.

Read \src{internal/guard/egress_facts.go}: destination, method and credential
position are facts the policy checks. \src{internal/guard/taint.go} reads
provenance; \src{internal/guard/textsecret.go} uses the shape detector in
\src{internal/detect/l1/l1.go}. These mechanisms answer different questions.
A value without a recognised key shape can still carry external provenance.
An approved destination does not make an unrelated credential appropriate.

For an MCP call continue through \src{internal/core/mcp.go},
\src{internal/mcp/client.go}, and either the HTTP or stdio transport.
Contained servers also involve \src{internal/core/contained.go},
\src{internal/container/server.go} and \src{internal/egress/proxy.go}.
Trace both the tool call and the server's onward network request; approving
one is not a general approval of the other. Study \file{test/exfil/}
alongside these files. Its Scheme scenarios at the end of the catalogue
are adversarial inputs, not setup commands to run in your working session.

\section{Route four: the same session through the console}

Start with the method and event shapes in
\src{internal/frontend/rpc/types.go} and
\src{internal/frontend/rpc/protocol.go}. Follow connection handling in
\src{internal/frontend/rpc/server.go} to dispatch in
\src{internal/frontend/rpc/sessions.go}, then the feature-specific handler:
sessions, workspaces, credentials, settings or runs. The daemon owns the
session; the browser owns a view of it.

Now read \src{internal/frontend/web/server.go}. Its \code{allowed} method
is a second boundary: a daemon method is not automatically available to a
browser, and a setup operation may be local-only. In the assets directory,
find the connection capability check, then the page action and its render.
If a new action appears but does nothing, inspect the request and the
advertised capabilities before changing the renderer. Go tests cannot
execute these modules: \file{test/browser/checks/} and
\file{test/browser/run.sh} exercise the actual browser.

\section{Route five: resuming without rerunning effects}

Pair \src{internal/core/resume.go} with \src{internal/log/payloads.go}.
Resuming assembles a runtime from today's configuration and init file;
it does not re-evaluate every old Scheme definition or rerun unfinished
tool calls. Compare this with \src{internal/window/replay.go}: replaying
a request reconstructs its recorded recipe and checks its hash without
asking a model or executing tools. Read
\file{internal/core/resume\_test.go} and
\file{internal/window/replay\_test.go} before altering either path.

\chapter{The repository customises its own agent}
\label{chap:repo-customisation}

This is a worked example from this checkout, not a hypothetical plugin.
Keep \src{.nodwyn/init.scm}, \src{dev/go.scm}, \src{dev/git.scm} and
\src{dev/repo.scm} open together. They contain the whole customisation;
the catalogue reproduces them without shortening the functions.

\section{Who writes the policy, and what loading means}

Anyone who can commit to a repository can write its init file, its helper
files and its \file{.nodwyn/allow.toml}, including the author of a clone
you have never inspected. They are project-supplied policy, not evidence
that a remote author deserves authority. ADR 45 requires an external user trust
record before either init or approvals load: review with \code{nodwyn trust},
then accept the exact displayed revision. A change to either file requires
review again; the checked bytes are the loaded bytes. The registry lives in
the user configuration directory, and other project dependencies are not hashed.
ADRs 14 and 17 describe the
project-init and standing-approval decisions. Read those decisions with
\src{internal/core/construct.go} and \src{internal/guard/allow.go}.

After trust, initialisation evaluates the init file with user origin, but through the
guard. While it loads, confirmations are declined because no frontend is
listening yet. User origin does not remove deny rules. Standing approvals
can waive a confirmation; they cannot turn a denial into an allow.
Protected policy paths cannot be included in ordinary write approvals.
Review a clone's policy before treating it as your working environment.

The first three forms load \file{dev/go.scm}, \file{dev/git.scm} and
\file{dev/repo.scm}, in that order. This matters: the latter files call
\code{run}, defined in the first. Loading registers functions and tools;
their bodies do not execute just because they were defined. The init file
then extends \code{look-tools} so its read-only navigation tools count
towards the default loop's reading reminder.

\section{Small tools over the existing boundary}

Read \code{run} in \src{dev/go.scm}. It calls \code{shell-exec}, then
formats status, stdout and stderr for the model. Its string is a report,
not an exception on every nonzero exit. \code{go-test} adds \code{-race};
\code{go-build}, \code{go-vet}, \code{go-fmt-check} and \code{run-make}
are named argument-list wrappers. The model sees descriptions and schemas
from \code{register-tool}; the guard still judges the inner shell call.

In \src{dev/git.scm}, \code{git-log} adds a read-only history tool.
There is no commit or push wrapper. In \src{dev/repo.scm}, compare
\code{grep-repo} and \code{find-files}, which use git, with
\code{read-adr}, which composes \code{fs-list} and \code{fs-read}.
The latter finds a zero-padded ADR prefix using \code{pad4}; an absent ADR
raises an ordinary Scheme error. These are two ways to reuse host
capabilities without adding a Go primitive.

The standing approvals printed after the init listing explain why these
commands can usually proceed without repeated questions. The write list
covers implementation, documentation, tests and \file{dev/}; it excludes
\file{.nodwyn/}, \file{scripts/} and the Makefile. This is why the validation
wrapper is defined in the protected init file, not in a freely writable
helper beside the other development tools.

\section{A validation wrapper, line by line in execution order}

\begin{enumerate}
\item \code{validation-state} combines git status with the tracked diff.
  Status alone would miss a second edit to an already changed tracked file.
  The source explicitly notes a limitation: editing an already untracked
  file again can leave both observations unchanged.
\item \code{validated-state} keeps its old binding if present. The
  \code{guard} expression supplies a default only when the variable is not
  defined yet. This makes reloading preserve the last observed state.
\item The \code{validate} tool starts \code{make validate} with
  \code{shell-start}. \code{validate-read} repeatedly uses
  \code{shell-read}, keeping output until the process has exited and no
  more buffered bytes remain. This avoids the synchronous shell timeout.
\item After collecting the report it stores the current repository state.
  Notice the exact promise: this records that validation \emph{ran}, not
  that it passed. It does not parse the final PASS/FAIL line. The agent must
  read failures and warnings, fix what belongs to the change, rerun, and
  report the result. A reminder is not a proof of correctness.
\item \code{validate-base-agent-turn} and
  \code{validate-base-test-call?} capture the original procedures once.
  Without the existing-binding check, reload would wrap the wrapper again.
  The new \code{test-call?} recognises \code{validate} as well as the
  ordinary test commands recognised by the base loop.
\item The replacement \code{default-agent-turn} snapshots the repository,
  runs the captured strategy, and compares the final state with both the
  starting and last-validated states. If work changed and is not validated,
  it adds a reminder and continues the conversation, at most twice.
  \code{validate-continue} runs calls, compacts and counts reading rounds
  in the same pattern as the normal chat loop.
\end{enumerate}

The companion implementation is \src{scripts/validate/main.go} and its
neighbouring files. Read \src{scripts/validate/changes.go} to understand the
comparison base, \src{scripts/validate/checks.go} for check selection,
\src{scripts/validate/base.go} for regression checks on unchanged code,
and \src{scripts/validate/runner.go} for process results. The Scheme
wrapper arranges when to ask for evidence; this Go program gathers it.

\section{Reload is a lifecycle boundary}

Read \src{internal/core/reload.go} before experimenting. Reload can
replace a strategy or the body of an existing tool for subsequent use.
Tool schemas, guard checks and fixed prompts need a new session or resume;
reload returns notes about changes that cannot take effect in place.
The model cannot call this primitive as a normal model-origin evaluation.

The repository example is not wholly idempotent: extending
\code{look-tools} with \code{append} again can add duplicate names, although
membership checks still work. The captured base procedures, by contrast,
must be preserved to avoid recursive wrapping. Ask which kind of state
each new definition represents before using this pattern elsewhere.

\chapter{Reading the Scheme layers in dependency order}

\section{The language below the strategy}

Start with \src{internal/scheme/value.go}, \src{internal/scheme/reader.go}
and \src{internal/scheme/env.go}: values, syntax and lexical bindings.
Then read \src{internal/scheme/interp.go}, moving to the
\file{forms\_*.go} files only when the evaluator dispatches to them.
\src{internal/scheme/prelude/base.scm} builds convenience procedures over
the Go builtins. Compare one definition, such as a list operation, with
the builtin it calls instead of trying to memorise the whole prelude.

\src{internal/scheme/prim.go} is the extension contract. Parameter types
make schemas; roles tell the guard what arguments mean. Follow a value's
taint through \src{internal/scheme/taint.go} and the taint builtins.
Ordinary errors, denials and halts have different meanings; read
\src{internal/scheme/errors.go} and \src{internal/scheme/halt.go} before
assuming that Scheme's \code{guard} can catch every failure.

\section{Default loop and its wrappers}

The following order follows \code{loop.New}. Later files may refer to
earlier definitions or capture them before redefining their public names.
Inspect the captured \code{*-base-*} bindings whenever a wrapper's behaviour
is surprising. Definition order and call order are different: defining a
procedure can refer to a global that will exist by the time it is called.

\begin{description}
\item[\src{internal/loop/strategy.scm}] Find \code{nodwyn-turn},
  \code{agent-turn}, \code{default-agent-turn} and \code{ask-model} first.
  Then read output-cut retries, reading reminders and \code{evidence-loop}.
  Separate the strategy's reminders from the host's enforceable stops.
\item[\src{internal/loop/prompts.scm}] Read the system prompt and the
  smaller prompt-producing procedures with their callers. These strings
  instruct a model; they do not replace host checks. Comments connect
  measured working rules to ADR 32.
\item[\src{internal/loop/compact.scm}] Follow threshold detection, the
  handover request and the context operation that uses the result. Pair it
  with \file{internal/core/compact\_test.go}; compaction changes a view of
  history, not the historical log.
\item[\src{internal/loop/agents.scm}] Follow the sub-agent's context and
  tool subset, then its result returning to the caller. Read the dispatch
  wrapper as well as \code{sub-agent}: agent tools participate in ordinary
  tool accounting. Pair it with \file{internal/core/agents\_test.go}.
\item[\src{internal/loop/optimize.scm}] Trace the specialised sub-agent
  which proposes improvements to project customisation. Read its limits
  and the protected write path before assuming it can silently change
  policy. Compare \file{internal/core/optimize\_test.go}.
\item[\src{internal/loop/plan.scm}] Follow \code{plan-gate} from strategy
  into the host's plan primitives in \src{internal/core/plan.go}. Scheme
  arranges planning; the host retains the irreversible-action gate.
\item[\src{internal/loop/shell.scm}] Read the innermost dispatch wrapper:
  how a shell-style command is recognised and redirected, and how results
  are folded. Compare \file{internal/core/habits\_test.go}.
\item[\src{internal/loop/hygiene.scm}] Track reads, edits and test outputs,
  then the conditions that make folding worthwhile. Its model-turn
  wrapper observes provider cache usage; its sub-agent wrapper distinguishes
  the main context from nested work.
\item[\src{internal/loop/readdef.scm}] Follow definition search into
  bounded file reads, and how those reads feed hygiene's bookkeeping.
  The model's navigation tool is Scheme composition over host primitives.
\item[\src{internal/loop/envhint.scm}] Follow the dispatch wrapper's
  detection of missing-module errors and the hint it appends to the folded
  result. Its per-session counter limits reminders to use the project's
  own environment. This is advice, not permission to execute a command.
\item[\src{internal/loop/checkloop.scm}] Read the outer turn wrapper and
  the independent-reading tools. A saved check must represent the user's
  actual example and survive compaction; the introductory comments explain
  the failures this rule addresses.
\end{description}

Finally read the project init and development helpers, followed by the
nineteen exfiltration scenarios. For each scenario identify the primitive,
the relevant argument role, the checker expected to act and the test fixture
which supplies its files or credentials. The scenario's filename is a hint;
\file{test/exfil/exfil\_test.go} supplies the assertions and setup.

\chapter{Finding the next change and proving it}

\section{Choose the owner before changing code}

A new composed workflow usually begins in Scheme. A new operation on the
host begins with a primitive and its metadata. A new provider dialect
belongs in \file{llm/}; provider-neutral context selection belongs in
\file{window/}. A new frontend action follows the RPC type, handler,
bridge capability and UI action together. A new persistent fact needs a
log payload and a replay/resume decision, not just an in-memory field.

Use the catalogue as a searchable symbol reference. Every file has its
own PDF bookmark and complete numbered listing. In a checkout, these
commands locate definitions and their consumers without a model call:
\begin{sh}
rg -n 'func .*Send|func .*startTurn' internal/core
rg -n 'applyPrim|Before\(' internal/scheme internal/guard
rg -n 'define .*run-calls|base-run-calls' internal/loop
rg -n 'KindToolResult|ToolResult' internal/log internal/loop internal/window
rg -n '^func Test' internal/core/strategy_test.go
\end{sh}

Read the surrounding test before modifying a contract. The fake provider
in \src{test/fakes/ollama/ollama.go} is scripted: \code{NewT} fails a test
if the loop asks after its script ends. The fake home helper in
\src{test/fakes/home/home.go} isolates tests which construct a guard or
commit. The catalogue includes the other fake services so readers can
follow integration tests without a paid model, real key or passphrase.

\section{Three study exercises with a way to check your answer}

\paragraph{Explain a refused write.} Pick one write scenario from the
catalogue. Draw the call from model arguments through metadata and the
guard to the uncalled handler. Verify the expected refusal in its Go test.
Your explanation is incomplete if it cannot name the argument carrying
the path and the argument carrying the payload.

\paragraph{Explain a context fold.} Follow one hygiene-fold operation
into the host and then into \code{window.Build}. Find where the old event
remains available and where its replacement appears in a later request.
Use \file{internal/core/habits\_test.go} and
\file{internal/core/strategy\_test.go} to check the distinction.

\paragraph{Explain the custom finish check.} On paper, simulate three
turns: no edit; an edit followed by validation; and an edit after validation.
Compare \code{before}, \code{now} and \code{validated-state} at each exit.
Then repeat with a failed validation and an already-untracked file. The
correct answer includes the two limitations described in the customisation
chapter rather than promising that the wrapper proves every change sound.

\section{Keeping the printed source honest}

\file{docs/book/source/generate.py} generates only the catalogue and its
coverage totals. Listings use \code{lstinputlisting} to read the actual
repository files at PDF build time; there is no second copy of the code
to edit. Both books include the same generated file. After adding, removing
or renaming source files, regenerate the catalogue and review its diff:
\begin{sh}
python3 docs/book/source/generate.py
python3 docs/book/source/generate.py --check
make book
make expert
make validate
\end{sh}

The generator enumerates git-tracked files. Add a new source file to git's
index before regenerating, or it will not yet be in the inventory. Its
check mode fails when the inventory or measured line counts are stale.
Source text edits are picked up by LaTeX directly; regenerate too if line
counts change. PDFs are a snapshot of the checkout used for that build,
not a claim that a particular release contains every listed function.

The repository's validation command selects checks by changed paths.
Changes confined to the shared source directory are under \file{docs/book/};
run \code{make expert} explicitly as well, because the validator's path
selection does not infer that shared dependency. When changing Go, follow
the package and whole-repository checks in \file{AGENTS.md}. When changing
the console, run its browser checks. A successful typeset is evidence for
the book build, not evidence that changed application code works.
